\documentclass[UTF8,11pt,a4paper]{ctexart}
\usepackage[a4paper,left=1.45cm,right=1.45cm,top=1.55cm,bottom=1.65cm]{geometry}
\usepackage{amsmath,amssymb,bm}
\usepackage{graphicx}
\usepackage{booktabs,tabularx,array,longtable}
\usepackage{enumitem}
\usepackage[dvipsnames,svgnames]{xcolor}
\usepackage{tikz}
\usetikzlibrary{arrows.meta,positioning,calc,fit,shapes.geometric}
\usepackage[most]{tcolorbox}
\usepackage{hyperref}
\usepackage{caption}
\usepackage{float}
\usepackage{fancyhdr}
\usepackage{titlesec}
\usepackage{microtype}
\IfFontExistsTF{Times New Roman}{\setmainfont{Times New Roman}}{\setmainfont{Liberation Serif}}
\IfFontExistsTF{Noto Sans CJK SC}{\setCJKmainfont{Noto Sans CJK SC}}{}
\IfFontExistsTF{Noto Sans CJK SC}{\setCJKsansfont{Noto Sans CJK SC}}{}
\IfFontExistsTF{DejaVu Sans Mono}{\setmonofont{DejaVu Sans Mono}}{\setmonofont{Latin Modern Mono}}
\definecolor{mainblue}{HTML}{0B5CAD}
\definecolor{softblue}{HTML}{EAF4FF}
\definecolor{darktext}{HTML}{1F2937}
\definecolor{softgray}{HTML}{F5F7FA}
\definecolor{accent}{HTML}{F59E0B}
\definecolor{goodgreen}{HTML}{2E7D32}
\definecolor{badred}{HTML}{B91C1C}
\definecolor{purple}{HTML}{6D5BD0}
\hypersetup{colorlinks=true,linkcolor=mainblue,urlcolor=DarkOrange,citecolor=mainblue}
\setlength{\parindent}{2em}
\setlength{\headheight}{14pt}
\setlength{\parskip}{0.22em}
\setlist[itemize]{leftmargin=2em,itemsep=0.23em,topsep=0.25em}
\setlist[enumerate]{leftmargin=2.2em,itemsep=0.23em,topsep=0.25em}
\renewcommand{\arraystretch}{1.2}
\captionsetup{font=small,labelfont=bf}
\titleformat{\section}{\Large\bfseries\sffamily\color{mainblue}}{\thesection}{0.7em}{}
\titleformat{\subsection}{\large\bfseries\sffamily\color{darktext}}{\thesubsection}{0.6em}{}
\titleformat{\subsubsection}{\normalsize\bfseries\sffamily\color{darktext}}{\thesubsubsection}{0.55em}{}
\pagestyle{fancy}
\fancyhf{}
\fancyhead[L]{\small\sffamily vHeat: Building Vision Models upon Heat Condu}
\fancyhead[R]{\small\sffamily \leftmark}
\fancyfoot[C]{\small\thepage}
\renewcommand{\headrulewidth}{0.4pt}
\newtcolorbox{infobox}[1]{enhanced,breakable,colback=softblue,colframe=mainblue!70!black,boxrule=0.8pt,arc=2mm,left=1.5mm,right=1.5mm,top=1mm,bottom=1mm,title=\bfseries #1}
\newtcolorbox{keybox}[1]{enhanced,breakable,colback=softgray,colframe=darktext!45,boxrule=0.6pt,arc=2mm,left=1.5mm,right=1.5mm,top=1mm,bottom=1mm,title=\bfseries #1}
\newtcolorbox{notebox}[1]{enhanced,breakable,colback=accent!8,colframe=accent!80!black,boxrule=0.7pt,arc=2mm,left=1.5mm,right=1.5mm,top=1mm,bottom=1mm,title=\bfseries #1}
\newcommand{\term}[1]{\textbf{\textcolor{mainblue}{#1}}}
\begin{document}

\section{论文基本信息}
\begin{center}
\begin{tabularx}{0.94\linewidth}{>{\bfseries}p{0.18\linewidth}X}
\toprule
题目 & vHeat: Building Vision Models upon Heat Conduction\\
课件日期 & 2025-05-22\\
研究方向 & 视觉模型 / 热传导 / 高效全局建模\\
一句话概括 & 用热传导方程类比视觉语义在二维空间中的扩散,以较低复杂度获得全局感受野,并构造物理启发的视觉骨干网络。\\
\bottomrule
\end{tabularx}
\end{center}

\begin{infobox}{摘要要点}
用热传导方程类比视觉语义在二维空间中的扩散,以较低复杂度获得全局感受野,并构造物理启发的视觉骨干网络。 本说明按“研究问题—核心机制—基础公式—实验阅读—局限与优化”的顺序重新组织，不保留原始材料的封面、目录和页码对应。
\end{infobox}

\section{研究问题与动机}
这项工作的出发点不是简单地替换某个网络层，而是识别现有范式中最影响\term{质量、效率、可靠性或可部署性}的瓶颈。结合论文所讨论的问题，可以将研究动机概括为以下几点：
\begin{itemize}
\item 引言 CNNs在视觉任务中的局限性(第一章) ViTs的高计算复杂度问题(第一章) 热传导原理与视觉语义传播的相似性(第一章) vHeat模型的提出背景和目标(第一章) Building Vision Models upon Heat Conduction
\item 相关工作 CNNs的发展历程及其局限性(第二章) ViTs的创新点及其计算复杂度问题(第二章) State Space Models在视觉任务中的应用(第二章) 生物学和物理学启发的模型研究现状(第二章) Building Vision Models upon Heat Conduction
\item 动态热扩散率 动态热扩散率的引入动机(第三章) 可学习的频率值嵌入(FVEs)的作用(第三章) 动态热扩散率对模型性能的影响(第三章) Building Vision Models upon Heat Conduction
\end{itemize}

\begin{notebox}{理解重点}
阅读这类论文时，应把“问题是否真实存在”“方法是否直接解决该问题”“实验是否排除了其他解释”分开判断。仅凭更大的模型、更多数据或更长训练得到的提升，不能自动视为结构创新带来的收益。
\end{notebox}

\section{核心思想与整体流程}
用热传导方程类比视觉语义在二维空间中的扩散,以较低复杂度获得全局感受野,并构造物理启发的视觉骨干网络。 从信息流角度看，其基本流程可以整理为下图。

\begin{figure}[H]
\centering
\begin{tikzpicture}[>=Latex,node distance=0.78cm]
\tikzstyle{procstep}=[draw,rounded corners,fill=softblue,align=center,minimum width=3.05cm,minimum height=0.95cm,text width=2.75cm]
\node[procstep] (a) {局部视觉特征};
\node[procstep,right=of a] (b) {热扩散算子};
\node[procstep,right=of b] (c) {频率调制};
\node[procstep,right=of c] (d) {全局表征};
\draw[->,thick,mainblue] (a)--(b);
\draw[->,thick,mainblue] (b)--(c);
\draw[->,thick,mainblue] (c)--(d);
\end{tikzpicture}
\caption{核心信息流与处理阶段。}
\end{figure}

\subsection{方法组成与信息流}
\begin{itemize}
\item 热传导原理 热传导方程的数学描述(第三章) 二维热传导方程的傅里叶变换解(第三章) 离散余弦变换(DCT)的实现方法(第三章) Building Vision Models upon Heat Conduction
\item vHeat模型架构 vHeat模型的整体架构(第三章) HCO层的设计细节(第三章) vHeat模型的Tiny/Small/Base三种规模(第三章) Building Vision Models upon Heat Conduction
\item 热传导算子(HCO) HCO的物理意义和数学实现(第三章) HCO的计算复杂度分析(第三章) HCO与自注意力机制的对比(第三章) Building Vision Models upon Heat Conduction
\item 引言 CNNs在视觉任务中的局限性(第一章) ViTs的高计算复杂度问题(第一章) 热传导原理与视觉语义传播的相似性(第一章) vHeat模型的提出背景和目标(第一章) Building Vision Models upon Heat Conduction
\item 相关工作 CNNs的发展历程及其局限性(第二章) ViTs的创新点及其计算复杂度问题(第二章) State Space Models在视觉任务中的应用(第二章) 生物学和物理学启发的模型研究现状(第二章) Building Vision Models upon Heat Conduction
\item vHeat模型通过引入热传导算子(HCO),实现了全局感受 野、低计算复杂度和高可解释性。(第五章) vHeat模型的潜在应用和未来发展方向。进一步应用于更多视觉任务,例如视频分析 和三维视觉。(第五章) Building Vision Models upon Heat Conduction
\item 模型总结与展望 vHeat模型的主要贡献。
\end{itemize}

\begin{keybox}{结构解析}
\begin{itemize}
\item 3 提出方法的优劣势 优势1： 全局感受野，能够捕捉图像中的长距离依赖关系（第三章） 优势2： 计算复杂度低（O(N\\textasciicircum{}1.5)），相比自注意力机制（O(N\\textasciicircum{}2)）效率更高（第三章） 优势3： 物理启发的可解释性，
\end{itemize}
\end{keybox}

上述内容以文字方式保留方法结构，避免把整张幻灯片作为独立页面嵌入正文。

\section{数学与机制理解}
本节给出理解该工作的基础数学结构。公式用于解释方法所属范式，并不替代原论文中的完整推导。
\begin{equation}
y(p)=\sum_{q\in\mathcal N(p)}K(p,q)\,x(q)
\end{equation}
视觉骨干的核心是把邻域或全局位置 $q$ 的信息聚合到位置 $p$。不同方法通过动态卷积、物理扩散、大核或多分支结构改变核 $K(p,q)$ 的范围与适应性。

\begin{keybox}{从公式回到方法}
应重点观察论文究竟改变了公式中的哪一部分：是输入表示、连接矩阵、条件变量、参数化方式、训练目标，还是求解与调度过程。真正的创新通常可以在这一层被准确定位。
\end{keybox}


\section{实验结果应如何阅读}
\begin{itemize}
\item 图像分类实验 ·vHeat在ImageNet-1K数据集上的性能(第四章) ·与其他模型的性能对比(第四章) ·vHeat的计算效率优势(第四章) Building Vision Models upon Heat Conduction
\item 语义分割实验 vHeat在ADE20K数据集上的语义分割性能(第四章) 与其他模型的性能对比(第四章) vHeat的计算效率优势(第四章) Building Vision Models upon Heat Conduction
\item 目标检测和实例分割实验 vHeat在COCO数据集上的目标检测和实例分割性能(第四章) 与其他模型的性能对比(第四章) vHeat的推理速度优势(第四章) Building Vision Models upon Heat Conduction
\item 低级视觉任务实验 vHeatIR在图像去噪和JPEG压缩伪影去除任务中的性能(第四章) 与其他模型的性能对比(第四章) Building Vision Models upon Heat Conduction
\item 动态热扩散率的消融实验 固定热扩散率与预测热扩散率的对比(第四章) 使用共享FVEs预测热扩散率的优势(第四章) Building Vision Models upon Heat Conduction
\item 鲁棒性评估 vHeat在ObjectNet和ImageNet-A数据集上的鲁棒性表现(第四章) 与其他模型的鲁棒性对比(第四章) Building Vision Models upon Heat Conduction
\item 热传导算子(HCO) HCO的物理意义和数学实现(第三章) HCO的计算复杂度分析(第三章) HCO与自注意力机制的对比(第三章) Building Vision Models upon Heat Conduction
\end{itemize}

\begin{notebox}{结果解读}
\begin{itemize}
\item 11 目标检测和实例分割实验 vHeat在COCO数据集上的目标检测和实例分割性能（第四章） 与其他模型的性能对比（第四章） vHeat的推理速度优势（第四章） Building Vision Models upon Heat Conduction
\end{itemize}
\end{notebox}

下图用于直观呈现主要对比、消融、定性样例或效率结果。
\begin{figure}[H]
\centering
\includegraphics[width=0.94\linewidth,height=0.43\textheight,keepaspectratio]{result_evidence.jpg}
\caption{代表性实验对比、消融结果或定性样例。}
\end{figure}
实验部分不应只看单一最好数字。还需要核对模型规模、训练数据、训练步数、采样或解码预算、硬件条件以及是否使用额外蒸馏、检索或后处理。只有比较条件接近时，结果才能真正说明方法本身的优势。

\section{优点、局限与可优化空间}
\subsection{主要优点}
\begin{itemize}
\item \term{问题与方法对应明确}：论文围绕一个可识别的核心瓶颈展开，方法结构能够直接映射到该瓶颈。
\item \term{可复用的设计思想}：即使不完整复现模型，其中的分解、稀疏化、递归、条件控制或系统优化思想也可迁移到相邻任务。
\item \term{实验提供了多角度证据}：实验通常同时展示主结果、效率比较、案例或消融，使结论不完全依赖单一指标。
\end{itemize}

\subsection{局限与进一步优化}
\begin{itemize}
\item 需要进一步区分\term{结构收益}与更大算力、更多数据、不同训练技巧带来的收益，并在等参数、等计算预算下进行严格比较。
\item 对超参数、输入分布和模型规模的敏感性仍应系统分析；在一个数据集上的最优配置不一定能直接迁移。
\item 实际部署还要考虑显存峰值、并发、延迟抖动、失败案例和鲁棒性，而不仅是离线平均指标。
\item 可尝试将该方法与蒸馏、量化、缓存复用、动态路由或更强的表示学习结合，验证不同优化是否互补。
\end{itemize}

\section{结论}
总体而言，用热传导方程类比视觉语义在二维空间中的扩散，以较低复杂度获得全局感受野，并构造物理启发的视觉骨干网络。。 进一步需要注意：vHeat模型通过引入热传导算子(HCO),实现了全局感受 野、低计算复杂度和高可解释性。(第五章) vHeat模型的潜在应用和未来发展方向。进一步应用于更多视觉任务,例如视频分析 和三维视觉。(第五章) Building Vision Models upon Heat Conduction;模型总结与展望 vHeat模型的主要贡献。。

\end{document}
